L'espace est muni d'un repère orthonormé \Oij.

On considère les points A$(1~;2~;~-1)$, B$(0~;3~;~2)$, C$(-2~;4~;~0)$ et D$(8~;2~;~-11)$.

\begin{enumerate}
\item
\begin{enumerate}[label=\textbf{\alph*.}]
\item 
Les points A, B et C définissent un plan si et seulement si ils ne sont pas alignés, si et seulement si les vecteurs $\vectt{AB}$ et $\vectt{AC}$ ne sont pas colinéaires.

Le vecteur $\vectt{AB}$ a pour coordonnées: 
$\begin{pmatrix}
x_{\text B} - x_{\text A} \\ y_{\text B} - y_{\text A} \\ z_{\text B} - z_{\text A}
\end{pmatrix}
=
\begin{pmatrix}
0 - 1 \\ 3-2 \\ 2-(-1) 
\end{pmatrix}
=
\begin{pmatrix}
-1 \\ 1 \\3
\end{pmatrix}$

De même le vecteur $\vectt{AC}$ a pour coordonnées: 
$\begin{pmatrix}
x_{\text C} - x_{\text A} \\ y_{\text C} - y_{\text A} \\ z_{\text C} - z_{\text A}
\end{pmatrix}
=
\begin{pmatrix}
-2-1 \\ 4-2 \\ 0-(-1)
\end{pmatrix}
=
\begin{pmatrix}
-3 \\ 2 \\ 1
\end{pmatrix}$

Or $x_{\vectt{AB}} \times 3 = x_{\vectt{AC}}$ et $y_{\vectt{AB}} \times 3 = 3 \neq y_{\vectt{AC}}$; donc les vecteurs $\vectt{AB}$ et $\vectt{AC}$ ne sont pas colinéaires.

On en conclut que les trois points A, B et C définissent un plan.  

\item
On considère le vecteur $\vect{n}\begin{pmatrix}5\\8\\-1\end{pmatrix}$.

$\vect{n} \cdot \vectt{AB} = 5 \times (-1) + 8 \times 1 +(-1) \times 3 = -5+8-3=0$ donc $\vect{n} \perp \vectt{AB}$.

$\vect{n} \cdot \vectt{AC} = 5 \times (-3) + 8 \times 2 +(-1) \times 1 = -15+16-1=0$ donc $\vect{n} \perp \vectt{AC}$.

Le vecteur $\vect{n}$ est donc orthogonal à deux vecteurs directeurs du plan (ABC) donc $\vect{n}$ est un vecteur normal au plan (ABC).

\item 
Le plan (ABC) est l'ensemble des points M\,$(x\;;\;y\;;\;z)$ tels que $\vectt{AM} \perp \vect{n}$.

$\vectt{AM}$ a pour coordonnées 
$\begin{pmatrix}
x-1 \\ y-2 \\ z+1
\end{pmatrix}$

$\begin{aligned}
\vectt{AM} \perp \vect{n}
\iff \vectt{AM} \cdot \vect{n}=0
& \iff (x-1) \times 5 + (y-2) \times 8 + (z+1) \times (-1)=0\\
& \iff 5x-5 +8y-16-z-1=0
\iff 5x+8y-z-22=0
\end{aligned}$

Donc le plan (ABC) a pour équation cartésienne: $5x+8y-z-22=0$.

\item 
On regarde si les coordonnées de D vérifient l'équation du plan.

$5x_{\text D}+8y_{\text D}-z_{\text D}-22
=8\times 5 + 8\times 2 - (-11) -22
= 40 +16 +11 -22 = 45 \neq 0$

Donc le point D n'appartient pas au plan (ABC).

\end{enumerate}

\item
	\begin{enumerate}[label=\textbf{\alph*.}]
		\item On considère la droite $\Delta$, orthogonale au plan (ABC) et passant par D.










La droite $\Delta$ est orthogonale au plan (ABC) donc elle a pour vecteur directeur n'importe quel vecteur normal au plan, en particulier le vecteur $\vect{n}$.

La droite $\Delta$  est donc l'ensemble des points M\,$(x\;;\;y\;;\;z)$ tels que $\vectt{DM}$ et $\vect{n}$ soient colinéaires, c'est-à-dire tels que $\vectt{DM}=t.\vect{n}$ avec $t\in\R$.

$\vectt{DM}=t.\vect{n}
\iff
\left \{
\begin{array}{l !{=} l}
x-8 & 5t\\
y-2 & 8t\\
z-(-11) & -t
\end{array}
\right .
\iff
\left \{
\begin{array}{l !{=} l}
x & 8+5t\\
y & 2+8t\\
z & -11 -t
\end{array}
\right .$

Donc une représentation paramétrique de la droite $\Delta$ est:
$\left \{
\begin{array}{l !{=} l}
x & 8+5t\\
y & 2+8t, \text{~~avec } t\in\R\\
z & -11 -t
\end{array}
\right .$

		\item On appelle E le point d'intersection de la droite $\Delta$ et du plan (ABC).



\begin{list}{\textbullet}{}
\item Le point E appartient au plan (ABC) donc $5x_{\text E}+8y_{\text E}-z_{\text E}-22=0$.
\item Le point E appartient à la droite $\Delta$ donc il existe un réel $t$ tel que
$\left \{
\begin{array}{l !{=} l}
x_{\text E} & 8+5t\\
y_{\text E} & 2+8t\\
z_{\text E} & -11 -t
\end{array}
\right .$
\end{list}

Le réel $t$ vérifie donc l'équation:
$5\left (8+5t\right ) + 8\left (2+8t\right ) - \left (-11-t\right )-22=0$,
c'est-à-dire:
$40+25t+16+64t+11+t-22=0$
soit:
$90t +45=0$
donc: $t=-\dfrac{1}{2}$

On a donc:
$\left \{
\begin{array}{l !{=} l}
x_{\text E} & 8+5 \times \left (- \frac{1}{2}\right )\\
y_{\text E} & 2+8\times \left (- \frac{1}{2}\right )\\
z_{\text E} & -11 - \left (- \frac{1}{2}\right )
\end{array}
\right .$
c'est-à-dire:
$\left \{
\begin{array}{l !{=} l}
x_{\text E} &  \frac{11}{2}\\
y_{\text E} & -2\\
z_{\text E} & - \frac{21}{2}
\end{array}
\right .$

		\item La valeur exacte de la longueur DE, distance du point D au plan (ABC) est: 
		
$\begin{aligned}
\text{DE}
& =\ds\sqrt{\left (x_{\text E} - x_{\text D}\right )^2 + \left (y_{\text E} - y_{\text D}\right )^2 + \left (z_{\text E} - z_{\text D}\right )^2}
 =\ds\sqrt{\left (\frac{11}{2} - 8\right )^2 + \left (-2 - 2\right )^2 + \left (-\frac{21}{2} - (-11)\right )^2}\\
 & = \ds\sqrt{\left (-\frac{5}{2}\right )^2 + \left (-4\right )^2 + \left (\frac{1}{2}\right )^2}
 = \ds\sqrt{\frac{25}{4} + 16 + \frac{1}{4} }
  = \ds\sqrt{\frac{90}{4}}
  = \dfrac{3\sqrt{10}}{2}
\end{aligned}$		
		
		  
	\end{enumerate}
\item On note H le point de coordonnées $\left(-\dfrac{2}{3}~;\dfrac{10}{3}~;\dfrac{4}{3}\right)$.

On admet qu'une représentation paramétrique de la droite (BC) est~:

\[
\begin{cases}
x = \phantom{3}-2k\\
y = 3+k\\
z = 2-2k
\end{cases}, \quad \text{avec } k \in \R.
\]

Avec un logiciel de calcul formel, on obtient les résultats suivants~:

\begin{center}
\begin{tabular}{|p{8.5cm}|}\hline
\textbf{Entrée 1~:} $(-2/3-1)*(-2)+(10/3-2)*1+(4/3-(-1))*(-2)$\\
\textbf{Sortie 1~:} $\rightarrow 0$\\
\hline
\textbf{Entrée 2~:} RésoudreSystème$(-2k=-2/3,\ 3+k=10/3,\ 2-2k=4/3)$\\
\textbf{Sortie 2~:} $\rightarrow k = 1/3$\\
\hline
\textbf{Entrée 3~:} RacineCarrée$\big((-2/3-1)^2+(10/3-2)^2+(4/3+1)^2\big)$\\
\textbf{Sortie 3~:} $\rightarrow$ RacineCarrée$(10)$\\
\hline
\textbf{Entrée 4~:} RacineCarrée$\big((-2-0)^2+(4-3)^2+(0-2)^2\big)$\\
\textbf{Sortie 4~:} $\rightarrow 3$\\
\hline
\end{tabular}
\end{center}

	\begin{enumerate}[label=\textbf{\alph*.}]
		\item
\begin{list}{\textbullet}{}
\item $\vectt{AH}$ a pour coordonnées:
$\begin{pmatrix}
-\frac{2}{3}-1 \\ \frac{10}{3} -2 \\ \frac{4}{3}-(-1)
\end{pmatrix}$
et $\vect{BC}$:
$\begin{pmatrix}
-2-0 \\ 4-3 \\ 0-2
\end{pmatrix}
=
\begin{pmatrix}
-2 \\ 1 \\ -2
\end{pmatrix}$

Le premier résultat de la feuille de calculs signifie que le produit scalaire $\vectt{AH}\cdot \vectt{BC}$ est nul, donc que (AH) $\perp$ (BC).

\item D'après le 2\ieme{} résultat de la feuille de calculs, le système
$\left \{
\begin{array}{l !{=} l}
-2k & -\frac{2}{3}\\
3+k & \frac{10}{3}\\
2-2k & \frac{4}{3}
\end{array}
\right .$
admet une solution unique; donc les coordonnées de H vérifient la représentation paramétrique de la droite (BC).

Donc le point H appartient à la droite (BC).
\end{list}		
		
De ces deux résultats, on peut déduire que le point H est le projeté orthogonal du point A sur la droite (BC).		
		
		\item On veut déterminer le volume du tétraèdre ABCD.



Le volume du tétraèdre ABCD est égal à un tiers du produit de l'aire de la base ABC par la hauteur $h$ relative à cette base.
 
\begin{list}{\textbullet}{}
\item La hauteur $h$ est la distance entre le point D et le plan (ABC), c'est-à-dire DE soit $\dfrac{3\sqrt{10}}{2}$.

\item D'après la question précédente, [AH] est la hauteur du triangle ABC relative au côté [BC]; donc l'aire du triangle ABC est: 
$\mathcal{A}= \dfrac{\text{AH} \times \text{BC}}{2}$.

D'après le 3\ieme{} résultat de la feuille de calculs, on a: $\text{AH}=\sqrt{10}$.

D'après le 4\ieme{} résultat de la feuille de calculs, on a: $\text{BC}=3$.

On en déduit que $\mathcal{A}= \dfrac{3\sqrt{10}}{2}$.
\end{list}

Le volume du tétraèdre ABCD est égal à:
$\dfrac{1}{3}\times \dfrac{3\sqrt{10}}{2}\times \dfrac{3\sqrt{10}}{2}
= \dfrac{15}{2}$.
	\end{enumerate}
\end{enumerate}
